\documentclass[12pt]{article}

\usepackage[margin=1in]{geometry}
\usepackage{enumitem}
\usepackage{amsmath}
\usepackage{amssymb}
\usepackage{array}
\usepackage{titlesec}

\title{Lecture Notes: Aristotle's \textit{Topics}, Book VI}
\author{}
\date{}

\begin{document}

\maketitle

\section*{Central Theme}

Aristotle's \textit{Topics}, Book VI gives commonplace rules for testing
definitions.

\[
\boxed{
\textit{Topics}, \text{ Book VI}
=
\text{commonplace rules for testing definitions}
}
\]

Book VI teaches how to ask whether a proposed account really states what a
thing is.

\[
\boxed{
\text{A definition is not merely a description; it is an account of essence.}
}
\]

A definition must be clear, non-redundant, non-circular, prior, intelligible,
properly placed in genus, properly divided by differentiae, peculiar to the
object, and expressive of essence. A definition is not a metaphor, not a list,
not a synonym, not a cause, not an effect, and not a description of the best
case. It is an account of what the thing is.

\section*{1. From Property to Definition}

Book V asked:

\[
\text{Is } P \text{ really a property of } S?
\]

Book VI asks:

\[
\text{Is this really the definition of } S?
\]

A property is convertible with the subject but does not state the essence.

A definition is convertible with the subject and does state the essence.

\subsection*{Whiteboard}

\[
\text{Property} = \text{convertible + non-essential}
\]

\[
\text{Definition} = \text{convertible + essential}
\]

\subsection*{Teaching Point}

A definition does not merely distinguish a thing. It says what the thing is.

\[
\boxed{
\text{Definition states essence.}
}
\]

\section*{2. The Five Tests of Definition}

Aristotle says that the discussion of definitions falls into five parts.

\[
\begin{array}{ll}
1. & \text{Does the expression apply to what the term applies to?} \\
2. & \text{Has the object been placed in the proper genus?} \\
3. & \text{Is the expression peculiar to the object?} \\
4. & \text{Does the expression state the essence?} \\
5. & \text{Has the object been defined incorrectly?}
\end{array}
\]

\subsection*{Whiteboard}

\[
\text{true of the subject}
\]

\[
\text{right genus}
\]

\[
\text{peculiar}
\]

\[
\text{essential}
\]

\[
\text{correctly stated}
\]

\subsection*{Teaching Point}

A definition must be true, proper, peculiar, essential, and well formed.

\[
\boxed{
\text{Definition is the strictest kind of predication.}
}
\]

\section*{3. Book VI Uses the Previous Books}

Aristotle says that some questions about definition are handled by rules already
given.

If the definition is not true of the subject, use the rules about accident.

\[
\text{truth} \rightarrow \text{accident rules}
\]

If the object is not placed in the proper genus, use the rules about genus.

\[
\text{genus} \rightarrow \text{genus rules}
\]

If the expression is not peculiar to the object, use the rules about property.

\[
\text{peculiarity} \rightarrow \text{property rules}
\]

\subsection*{Teaching Point}

Definition is the most demanding predicable because it includes the others.

\[
\boxed{
\text{A definition must do the work of truth, genus, property, and essence together.}
}
\]

\section*{4. Incorrect Definition: Obscurity and Redundancy}

Aristotle first asks whether the definition has been given incorrectly.

Incorrectness has two major forms.

\[
\begin{array}{ll}
1. & \text{obscure language} \\
2. & \text{unnecessary length or redundancy}
\end{array}
\]

A definition should be clear because its purpose is to make the thing known.

A definition should avoid unnecessary additions because it should state exactly
what is needed to express the essence.

\subsection*{Teaching Point}

A definition fails if it does not clarify or if it says more than is needed.

\[
\boxed{
\text{Bad definition is either obscure or redundant.}
}
\]

\section*{5. Obscurity Through Ambiguity}

A definition is obscure if it uses ambiguous terms.

\subsection*{Examples}

\[
\text{Becoming is a passage into being.}
\]

\[
\text{Health is a balance of hot and cold elements.}
\]

The problem is that terms such as ``passage'' and ``balance'' can be understood
in more than one sense.

\subsection*{Teaching Point}

A definition must not hide behind equivocal language.

\[
\boxed{
\text{Ambiguous term} \Rightarrow \text{unclear definition}
}
\]

\section*{6. Obscurity in the Term Being Defined}

The term being defined may itself be ambiguous.

If the term has many senses, the definer must say which sense is being defined.

\[
\text{ambiguous subject} \Rightarrow \text{ambiguous definition}
\]

\subsection*{Teaching Point}

You cannot define clearly unless you first know which meaning is under
discussion.

\[
\boxed{
\text{Before defining, distinguish the senses.}
}
\]

\section*{7. Metaphorical Definitions}

A definition is bad if it uses metaphorical language.

\subsection*{Examples}

\[
\text{knowledge is unsupplantable}
\]

\[
\text{earth is a nurse}
\]

\[
\text{temperance is a harmony}
\]

A metaphor may suggest something, but it does not literally state the essence.

\subsection*{Teaching Point}

A definition must classify literally, not poetically.

\[
\boxed{
\text{Metaphor} \neq \text{definition}
}
\]

\section*{8. Unfamiliar or Strange Language}

Definitions also fail when they use unfamiliar or strange terms.

\subsection*{Examples}

\[
\text{eye} = \text{brow-shaded}
\]

\[
\text{marrow} = \text{bone-formed}
\]

Such phrases require explanation before they can explain anything else.

\subsection*{Teaching Point}

A definition should not require a second definition before it can be understood.

\[
\boxed{
\text{Unfamiliar language defeats the purpose of definition.}
}
\]

\section*{9. Worse Than Metaphor}

Some expressions are worse than metaphor because they are neither literal nor
genuinely metaphorical.

\subsection*{Example}

\[
\text{law} = \text{measure or image of what is naturally just}
\]

If the expression is literal, it is false. If it is metaphorical, the likeness
must be clear. If there is no likeness and no literal truth, the expression
clarifies nothing.

\subsection*{Whiteboard}

\[
\text{literal} \rightarrow \text{acceptable if true}
\]

\[
\text{metaphorical} \rightarrow \text{partly intelligible by likeness}
\]

\[
\text{neither literal nor metaphorical} \rightarrow \text{mere obscurity}
\]

\subsection*{Teaching Point}

A definition must not pretend to explain by a phrase that has no clear relation
to the thing.

\[
\boxed{
\text{A definition must make contact with the essence.}
}
\]

\section*{10. A Good Definition Should Reveal Its Contrary}

A correctly rendered definition should help indicate the definition of the
contrary.

\subsection*{Example}

If we understand the definition of useful, we should be able to understand the
definition of injurious.

\[
\text{definition of } A
\Rightarrow
\text{some clarity about contrary of } A
\]

\subsection*{Teaching Point}

Definitions sit within networks of opposition.

\[
\boxed{
\text{A definition that leaves the contrary unintelligible is suspect.}
}
\]

\section*{11. A Definition Should Identify What It Defines}

If a definition is stated by itself and it is not evident what it defines, it is
badly rendered.

Aristotle compares this to old paintings that needed inscriptions before anyone
could recognize the figures.

\subsection*{Teaching Point}

A definition should not need a label to tell us what it defines.

\[
\boxed{
\text{A definition should identify its subject by its content.}
}
\]

\section*{12. Redundancy: Universal Attributes}

A definition is redundant if it includes something that belongs universally to
everything or to everything under the same genus.

The genus should separate the object from things in general.

The differentia should separate it from things within the genus.

A term that belongs to everything separates nothing.

\[
\text{universal attribute} \Rightarrow \text{no distinction}
\]

\subsection*{Teaching Point}

Every part of a definition should contribute to distinguishing the essence.

\[
\boxed{
\text{A defining term must divide.}
}
\]

\section*{13. Redundancy: Removable Additions}

If something can be removed and the remaining expression still states the
essence and is peculiar to the subject, the removed part was superfluous.

\subsection*{Whiteboard}

\[
\text{definition} - \text{term}
\]

If the remainder is still:

\[
\text{peculiar + essential}
\]

then:

\[
\text{term was redundant.}
\]

\subsection*{Teaching Point}

A definition should contain no idle words.

\[
\boxed{
\text{If removal does not harm the definition, the removed part was unnecessary.}
}
\]

\section*{14. False Additions Are Worse Than Redundant Additions}

If part of the definition does not apply to everything in the species, the
definition is worse than merely redundant. It becomes false or non-convertible.

\subsection*{Example}

\[
\text{man} = \text{walking biped animal six feet high}
\]

Not every man is six feet high.

Therefore, the definition fails.

\[
\text{attribute not true of every } S
\Rightarrow
\text{definition not convertible}
\]

\subsection*{Teaching Point}

A definition must apply to every instance of the thing defined.

\[
\boxed{
\text{A definition cannot include an accidental feature not shared by the whole species.}
}
\]

\section*{15. Repetition: Saying the Same Thing Twice}

Aristotle warns against repeating the same content in a definition.

\subsection*{Example}

\[
\text{desire} = \text{conation for the pleasant}
\]

This risks repeating ``for the pleasant'' if desire already means conation
toward the pleasant.

A clearer example is saying that wisdom both defines and contemplates reality,
if definition is already a kind of contemplation.

\subsection*{Teaching Point}

Redundancy is not merely repeated words, but repeated meaning.

\[
\boxed{
\text{Do not predicate the same content twice.}
}
\]

\section*{16. Universal Plus Particular}

Another redundancy occurs when a universal is mentioned and then a particular
case already included under it is added.

\subsection*{Example}

\[
\text{Equity is a remission of what is expedient and just.}
\]

If the just is already a branch of the expedient, the addition is redundant.

\subsection*{Whiteboard}

\[
\text{universal} + \text{included particular} = \text{repetition}
\]

\subsection*{Teaching Point}

Do not add a part after already naming the whole.

\[
\boxed{
\text{Do not include what is already included.}
}
\]

\section*{17. Definition Must Use Prior and More Intelligible Terms}

The central positive rule is that a definition should be made through terms that
are prior and more intelligible.

The purpose of definition is to make the object known.

We make things known through what is prior and clearer.

\[
\text{definition} \rightarrow \text{prior and more intelligible terms}
\]

\subsection*{Teaching Point}

A definition explains by what is prior, not by what is posterior.

\[
\boxed{
\text{Definition should proceed from what is prior and clearer.}
}
\]

\section*{18. More Intelligible Absolutely and More Intelligible to Us}

Aristotle distinguishes two kinds of intelligibility.

\[
\text{more intelligible absolutely}
\]

and:

\[
\text{more intelligible to us}
\]

A point is prior to a line, a line to a plane, and a plane to a solid.

But to us, a solid may be easier to grasp because it is more available to
perception.

\subsection*{Whiteboard}

\[
\text{absolutely: point} \rightarrow \text{line} \rightarrow \text{plane} \rightarrow \text{solid}
\]

\[
\text{to us: solid} \rightarrow \text{plane} \rightarrow \text{line} \rightarrow \text{point}
\]

\subsection*{Teaching Point}

Scientific definition should use what is prior absolutely, even if teaching
sometimes begins with what is clearer to us.

\[
\boxed{
\text{Pedagogical order and scientific order are not always the same.}
}
\]

\section*{19. Genus and Differentiae Are Prior to Species}

A correct definition should proceed through genus and differentiae.

\[
\text{definition} = \text{genus} + \text{differentiae}
\]

The genus and differentiae are prior to the species because if they are removed,
the species is removed.

\[
\text{remove genus and differentia}
\Rightarrow
\text{remove species}
\]

\subsection*{Teaching Point}

A species is understood through the genus and differentiae that make it what it
is.

\[
\boxed{
\text{Genus and differentia explain the species.}
}
\]

\section*{20. Do Not Define the Prior Through the Posterior}

A definition fails when it explains a prior thing through a posterior thing.

\subsection*{Geometry Examples}

A point may be described as the limit of a line.

A line may be described as the limit of a plane.

A plane may be described as the limit of a solid.

These may be pedagogically useful, but scientifically they explain the prior
through the posterior.

\[
\text{posterior} \not\rightarrow \text{scientific definition of prior}
\]

\subsection*{Teaching Point}

A teaching shortcut is not always a scientific definition.

\[
\boxed{
\text{Scientific definition must respect order of being.}
}
\]

\section*{21. Do Not Define the Definite Through the Indefinite}

Aristotle says one should not define what is at rest and definite through what
is indefinite and in motion.

\[
\text{definite} > \text{indefinite}
\]

\[
\text{stable} > \text{moving}
\]

\subsection*{Teaching Point}

The clearer and more determinate should not be explained by the less
determinate.

\[
\boxed{
\text{Definition should move toward determinacy.}
}
\]

\section*{22. Do Not Define Opposites Through Opposites Except Where Necessary}

A common error is defining one opposite through another.

\[
\text{good} = \text{opposite of evil}
\]

Opposites are simultaneous by nature, so one is not naturally prior to the
other.

\subsection*{Exception}

Some terms are essentially relative.

\[
\text{double} \leftrightarrow \text{half}
\]

In such cases, one cannot understand one without the other.

\subsection*{Teaching Point}

Do not define through an opposite unless the thing is essentially relative.

\[
\boxed{
\text{Opposites are usually simultaneous, not explanatory priors.}
}
\]

\section*{23. Do Not Use the Term Defined in the Definition}

A definition is circular if it uses the term being defined.

This can happen secretly.

\subsection*{Example}

\[
\text{sun} = \text{star that appears by day}
\]

But if:

\[
\text{day} = \text{the passage of the sun over the earth}
\]

then the definition secretly includes sun.

\[
\text{sun} \rightarrow \text{day} \rightarrow \text{sun}
\]

\subsection*{Teaching Point}

A circular definition hides the thing defined inside the defining expression.

\[
\boxed{
\text{Do not define a thing through itself, even indirectly.}
}
\]

\section*{24. Do Not Define One Co-Ordinate by Another}

Co-ordinate members of the same division are simultaneous by nature.

\subsection*{Example}

Odd and even are co-ordinate differentiae of number.

\[
\text{odd}
\]

and:

\[
\text{even}
\]

Therefore, one should not define odd as:

\[
\text{that which is greater by one than an even number}
\]

\subsection*{Teaching Point}

Co-ordinates should not be used as though one were prior to the other.

\[
\boxed{
\text{Members of the same division do not define one another.}
}
\]

\section*{25. Do Not Define the Superior Through the Subordinate}

A definition fails when it defines a superior term through a subordinate one.

\subsection*{Example}

\[
\text{good} = \text{state of virtue}
\]

But virtue is itself a kind of good.

\[
\text{virtue} \subset \text{good}
\]

\subsection*{Teaching Point}

The wider and prior should not be defined through the narrower and posterior.

\[
\boxed{
\text{Do not define the superior through the subordinate.}
}
\]

\section*{26. The Definition Must Place the Object in a Genus}

A definition fails if it does not state what kind of thing the object is.

\subsection*{Bad Examples}

\[
\text{body} = \text{that which has three dimensions}
\]

\[
\text{man} = \text{that which knows how to count}
\]

These do not first say what the thing is.

\subsection*{Teaching Point}

The genus gives the first answer to the question: What is it?

\[
\boxed{
\text{Definition begins with genus.}
}
\]

\section*{27. Use the Proper and Nearest Genus}

It is not enough to give some broad genus. A definition should use the proper or
nearest genus.

If one uses a higher genus, one must add the differentiae that specify the
nearer genus.

\[
\text{nearest genus} > \text{remote genus alone}
\]

\subsection*{Example}

\[
\text{plant}
\]

does not specify:

\[
\text{tree}
\]

\subsection*{Teaching Point}

The nearer genus gives a sharper account of the essence.

\[
\boxed{
\text{Use the nearest genus when possible.}
}
\]

\section*{28. Relative Terms Must Be Defined in Relation to the Right Object}

If the term being defined is relative, its definition must mention the
correlate.

\subsection*{Example}

Knowledge should not merely be defined as:

\[
\text{incontrovertible conception}
\]

It must be defined in relation to what is knowable.

\[
\text{knowledge} \rightarrow \text{knowable}
\]

\subsection*{Teaching Point}

The essence of a relative is its being-toward another.

\[
\boxed{
\text{Relative terms require their correlates.}
}
\]

\section*{29. Define Relative Terms by the Better or Final Correlate}

When a term relates to many things, define it in relation to the better or
proper end.

\subsection*{Example}

Medicine may relate accidentally to disease, but essentially it aims at health.

\[
\text{medicine} \rightarrow \text{health}
\]

not primarily:

\[
\text{medicine} \rightarrow \text{disease}
\]

\subsection*{Teaching Point}

A capacity or science should be defined by its proper end, not by accidental
effects.

\[
\boxed{
\text{Define a power by what it is for.}
}
\]

\section*{30. Differentiae Must Belong to the Genus}

A definition must use differentiae that actually divide the stated genus.

\[
\text{genus} + \text{proper differentia} = \text{species}
\]

If the alleged differentia does not belong to the genus, it cannot define a
species of that genus.

\subsection*{Teaching Point}

Differentiae must be native to the genus they divide.

\[
\boxed{
\text{A foreign differentia cannot define a species.}
}
\]

\section*{31. Differentiae Must Have Co-Ordinate Alternatives}

A true differentia belongs within a division.

\subsection*{Example}

Animal may be divided by:

\[
\text{walking}
\]

\[
\text{flying}
\]

\[
\text{aquatic}
\]

\[
\text{biped}
\]

If the alleged differentia has no co-ordinate member in a division, it is
suspect.

\subsection*{Teaching Point}

A differentia is one branch within a real division of a genus.

\[
\boxed{
\text{Differentiae come in co-ordinate divisions.}
}
\]

\section*{32. The Differentia Added to the Genus Must Make a Species}

A specific differentia, when added to the genus, should make a species.

\[
\text{genus} + \text{differentia} = \text{species}
\]

If the addition does not produce a species, the alleged differentia is not a
true specific differentia.

\subsection*{Teaching Point}

A differentia must specify, not merely decorate.

\[
\boxed{
\text{A real differentia forms a species.}
}
\]

\section*{33. Beware Dividing a Genus by Negation}

Aristotle criticizes definitions like:

\[
\text{line} = \text{length without breadth}
\]

This divides length by a negation.

\[
\text{length with breadth}
\]

\[
\text{length without breadth}
\]

If length is the genus, then length would have to partake of one of its own
species.

\subsection*{Qualification}

Negation is allowed in definitions of privations.

\subsection*{Example}

Blindness must be defined by reference to not seeing where sight is naturally
due.

\subsection*{Teaching Point}

A genus should not be divided merely by negation unless the subject is truly
privative.

\[
\boxed{
\text{Negation is dangerous unless the object is a privation.}
}
\]

\section*{34. Do Not Render Species as Differentia}

A species should not be used as though it were a differentia.

\subsection*{Example}

\[
\text{contumely} = \text{insolence accompanied by jeering}
\]

But jeering may be a species of insolence, not a differentia of it.

\subsection*{Teaching Point}

A differentia divides a genus; it is not one of the species divided off.

\[
\boxed{
\text{Species} \neq \text{differentia}
}
\]

\section*{35. Do Not Render Genus as Differentia}

Sometimes a genus is mistakenly treated as a differentia.

\subsection*{Example}

\[
\text{virtue} = \text{good state}
\]

Here:

\[
\text{genus} = \text{state}
\]

and:

\[
\text{differentia} = \text{good}
\]

if ``good'' indicates what sort of state virtue is.

\subsection*{Teaching Point}

The genus states what the thing is; the differentia states what sort of that
thing it is.

\[
\boxed{
\text{Genus and differentia must not be reversed.}
}
\]

\section*{36. Differentia Should Signify a Quality}

The differentia should indicate a quality of the genus, not an individual thing.

\[
\text{differentia} \rightarrow \text{quality of genus}
\]

\subsection*{Teaching Point}

A differentia qualifies the genus.

\[
\boxed{
\text{Differentia tells what sort of genus-member the species is.}
}
\]

\section*{37. Differentia Is Not Accidental}

A differentia is not an accident.

It cannot both belong and not belong to the species.

\[
\text{differentia} \neq \text{accident}
\]

\[
\text{differentia belongs essentially}
\]

\subsection*{Teaching Point}

The differentia is part of what the thing is, not something that happens to it.

\[
\boxed{
\text{Essential difference is not accidental attribute.}
}
\]

\section*{38. Genus Is Not Predicated of the Differentia}

The genus is predicated of the species, not of the differentia itself.

\subsection*{Example}

Animal is predicated of man, ox, and other animals.

But animal is not predicated of the differentia itself, such as walking.

\subsection*{Teaching Point}

The genus belongs to the species, not to the differentia taken by itself.

\[
\boxed{
\text{Genus is predicated of species, not of differentiae as such.}
}
\]

\section*{39. Differentia Must Be Prior to the Species}

The differentia is posterior to the genus but prior to the species.

\[
\text{genus} \rightarrow \text{differentia} \rightarrow \text{species}
\]

\subsection*{Teaching Point}

The species is formed by the genus as qualified by differentia.

\[
\boxed{
\text{Differentia stands between genus and species.}
}
\]

\section*{40. Beware Differentiae from Unrelated Genera}

A differentia should not belong to a different genus unrelated to the genus
being defined.

If the same differentia is used across genera, the genera must be related or
fall under a higher common genus.

\[
D \text{ used across } G_1 \text{ and } G_2
\Rightarrow
G_1, G_2 \text{ need ordering or common genus}
\]

\subsection*{Teaching Point}

Differentiae must not drag the species into unrelated genera.

\[
\boxed{
\text{A differentia should fit the genus it divides.}
}
\]

\section*{41. Locality Is Usually Not Essence}

Aristotle warns against treating ``existence in'' something as a differentia of
essence.

Location usually does not differentiate one essence from another.

\subsection*{Qualification}

Terms like aquatic may not simply mean ``in water.'' They may indicate a
quality or natural mode of life.

\[
\text{in water} \neq \text{aquatic essence}
\]

\subsection*{Teaching Point}

A real differentia should express what sort of thing it is, not merely where it
is.

\[
\boxed{
\text{Place is usually not essence.}
}
\]

\section*{42. Affection Is Not Differentia}

An affection should not be used as a differentia.

An affection, when intensified, can destroy the subject.

A differentia preserves the thing's essence.

\[
\text{affection intensified} \rightarrow \text{may destroy}
\]

\[
\text{differentia} \rightarrow \text{preserves essence}
\]

\subsection*{Teaching Point}

A differentia is constitutive, not destructive.

\[
\boxed{
\text{Do not define essence by a destructive affection.}
}
\]

\section*{43. Define by Natural Function}

When defining relative terms or instruments, one should define them according to
their natural or proper function.

\subsection*{Example}

A strigil can be used to dip water, but that is not its proper function.

\[
\text{proper function} \neq \text{possible accidental use}
\]

\subsection*{Teaching Point}

A thing should be defined by what the prudent user or proper science uses it
for.

\[
\boxed{
\text{Define by natural function, not accidental use.}
}
\]

\section*{44. Use the Primary Relation}

If a term can be related to several things, define it by its primary relation.

\subsection*{Example}

Wisdom should be defined primarily as the virtue of the reasoning faculty, not
merely as the virtue of man or the soul.

Man and soul are wise through the reasoning faculty.

\[
\text{wisdom} \rightarrow \text{reasoning faculty}
\]

\subsection*{Teaching Point}

Definition should name the primary seat or relation, not a derivative one.

\[
\boxed{
\text{Define by what is primary.}
}
\]

\section*{45. Dispositions and Affections Must Belong to the Right Subject}

If something is defined as a disposition or affection, it must be located in the
thing that can receive that disposition or affection.

\subsection*{Bad Examples}

\[
\text{sleep} = \text{failure of sensation}
\]

\[
\text{pain} = \text{violent disruption of naturally conjoined parts}
\]

Such definitions may confuse the state with its cause or place the affection in
the wrong subject.

\subsection*{Teaching Point}

A definition must place states, affections, causes, and effects in the right
subject.

\[
\boxed{
\text{Locate the affection where it naturally belongs.}
}
\]

\section*{46. Do Not Confuse Cause and Effect}

Definitions can fail by confusing cause and effect.

A disruption of parts may cause pain, but it is not pain.

Failure of sensation may accompany or cause sleep, but it is not sleep itself.

\[
\text{cause} \neq \text{effect}
\]

\subsection*{Teaching Point}

A definition must state what the thing is, not merely what causes it or follows
from it.

\[
\boxed{
\text{Cause and essence are not automatically the same.}
}
\]

\section*{47. Watch Temporal Mismatch}

Definitions fail if the expression applies at one time but the term defined does
not.

\subsection*{Example}

\[
\text{immortal} = \text{living thing immune at present from destruction}
\]

This is ambiguous.

It might mean merely:

\[
\text{not destroyed now}
\]

which is too weak.

Or it might mean:

\[
\text{now such that it can never be destroyed}
\]

which is closer to immortal.

\subsection*{Teaching Point}

A definition must match the temporal force of the term defined.

\[
\boxed{
\text{Do not weaken an eternal term into a present-tense condition.}
}
\]

\section*{48. Define by Choice, Not Mere Ability, Where Appropriate}

Aristotle criticizes defining justice as an ability to distribute what is equal.

Justice concerns the person who chooses justly, not merely the person who has
the ability to distribute.

\[
\text{just person} \neq \text{merely able to distribute equally}
\]

\subsection*{Teaching Point}

Moral definitions must often refer to choice, not bare capacity.

\[
\boxed{
\text{Virtue is not mere ability.}
}
\]

\section*{49. Definition and Degree Must Match}

If the thing defined admits of more and less, the definition should also admit
of more and less.

If the definition admits degrees but the term does not, or vice versa, something
is wrong.

\[
\text{term admits degree} \Leftrightarrow \text{definition admits degree}
\]

\subsection*{Teaching Point}

The definition should behave like the thing defined.

\[
\boxed{
\text{Name and definition should vary together.}
}
\]

\section*{50. Compare Application in Two Cases}

If the term applies more to one case while the definition applies more to
another, the definition fails.

\subsection*{Example}

Fire may denote flame more than light.

But:

\[
\text{body consisting of the most rarefied particles}
\]

may fit light more than flame.

If so, the proposed definition does not track the term.

\subsection*{Whiteboard}

\[
\text{name more applies to } A
\]

\[
\text{definition more applies to } B
\]

\[
\therefore \text{definition fails}
\]

\subsection*{Teaching Point}

The definition and the name must apply most strongly to the same cases.

\[
\boxed{
\text{A definition must track the ordinary application of the term.}
}
\]

\section*{51. Beware Definitions Relative to Two Separate Things}

A definition can fail if it defines something by relation to two things
separately.

\subsection*{Example}

\[
\text{beautiful} = \text{pleasant to the eyes or to the ears}
\]

Then something pleasant to the eyes but not to the ears may become both
beautiful and not beautiful.

\subsection*{Teaching Point}

Disjunctive definitions can produce contradictions if the relations are not
unified.

\[
\boxed{
\text{Do not define one essence by disconnected alternatives.}
}
\]

\section*{52. Substitute Definitions for the Terms Inside the Definition}

Aristotle recommends defining the parts of the definition and checking whether
the resulting account still works.

\[
\text{definition of term}
\rightarrow
\text{substitute definitions of its parts}
\]

\subsection*{Teaching Point}

A definition should remain coherent when its own terms are unpacked.

\[
\boxed{
\text{Unpack the definition to test it.}
}
\]

\section*{53. Relative Definitions Must Mention the End}

If a relative term is defined, it should be defined in relation to its end or
proper object.

\subsection*{Example}

Desire should be defined not merely by the pleasant thing, but by pleasure, if
pleasure is the proper end.

\subsection*{Teaching Point}

Relative definitions should point toward the proper end.

\[
\boxed{
\text{Define relative terms by their proper object or end.}
}
\]

\section*{54. Do Not Define by Process When the End Is the Completion}

If the correlate is a process or activity, we must ask whether the end is really
the process itself or the completion of the process.

\[
\text{process} \neq \text{completion}
\]

\subsection*{Teaching Point}

A definition should identify whether the activity or its completion is the true
end.

\[
\boxed{
\text{Do not confuse activity with its completed end.}
}
\]

\section*{55. Add Quantity, Quality, Place, and Cause When Needed}

A definition may fail because it omits necessary differentiae such as quantity,
quality, place, or cause.

\subsection*{Examples}

\[
\text{night} = \text{shadow on the earth}
\]

\[
\text{earthquake} = \text{movement of the earth}
\]

\[
\text{wind} = \text{movement of air}
\]

These are too broad unless one specifies the relevant manner, amount, place, or
cause.

\subsection*{Teaching Point}

A definition must include all differentiae needed to capture the essence.

\[
\boxed{
\text{Do not omit the qualifying features that make the thing this thing.}
}
\]

\section*{56. Conations Must Often Be Defined by the Apparent Good}

In cases of desire or wishing, one must often include the word ``apparent.''

People do not always desire what is truly good or pleasant. They desire what
appears good or pleasant to them.

\[
\text{wishing} \rightarrow \text{apparent good}
\]

\[
\text{desire} \rightarrow \text{apparent pleasant}
\]

\subsection*{Teaching Point}

Psychological definitions must account for appearance, not only reality.

\[
\boxed{
\text{Desire follows the apparent good or pleasant.}
}
\]

\section*{57. Definitions of States Clarify Related Things}

If one defines a state, one also indirectly clarifies what has the state, what
lacks it, and the corresponding activity.

\subsection*{Example}

Defining knowledge also clarifies:

\[
\text{knower}
\]

\[
\text{ignorance}
\]

\[
\text{knowing}
\]

\[
\text{being ignorant}
\]

\subsection*{Teaching Point}

A definition radiates through a family of related terms.

\[
\boxed{
\text{Define one state, and its correlative terms become clearer.}
}
\]

\section*{58. Relative Species Should Correspond to Relative Species}

If a genus is relative to something, then a species under that genus should be
relative to a species of that correlate.

\subsection*{Example}

If belief is relative to an object of belief, then a particular belief should be
relative to a particular object of belief.

\[
\text{genus relative to correlate}
\Rightarrow
\text{species relative to species of correlate}
\]

\subsection*{Teaching Point}

Relative structure should descend from genus to species.

\[
\boxed{
\text{Relations must remain ordered through the hierarchy.}
}
\]

\section*{59. Opposite Definitions Should Correspond}

If one term is defined, its opposite should have the opposite definition.

\subsection*{Example}

If double is what exceeds another by an equal amount to that other, then half is
what is exceeded by an amount equal to itself.

\[
\text{double} \leftrightarrow \text{half}
\]

\[
\text{definition of double} \leftrightarrow \text{definition of half}
\]

\subsection*{Teaching Point}

Correct definitions preserve the structure of opposition.

\[
\boxed{
\text{Opposite terms should have opposite accounts.}
}
\]

\section*{60. Be Careful with Privation}

A privative term must be defined by naming:

\[
\text{what is absent}
\]

and:

\[
\text{what is deprived}
\]

\subsection*{Example}

\[
\text{blindness} = \text{privation of sight in an eye}
\]

It is not enough to say ``privation'' generally.

\subsection*{Teaching Point}

A privation must say both what is lacking and where it is naturally due.

\[
\boxed{
\text{Privation requires subject and absent state.}
}
\]

\section*{61. Do Not Define as Privation What Is Not Mere Privation}

Not every negative condition is a privation.

\subsection*{Example}

Error is not simply absence of knowledge.

Error belongs to someone deceived, not merely to something without knowledge.

We do not call stones or children ``in error'' simply because they lack
knowledge.

\[
\text{error} \neq \text{mere lack of knowledge}
\]

\subsection*{Teaching Point}

Not every negative condition is a privation.

\[
\boxed{
\text{A privation is a structured lack, not mere absence.}
}
\]

\section*{62. Inflexions Should Match}

If the definition of a term is correct, its inflected forms should correspond.

\subsection*{Example}

If:

\[
\text{beneficial} = \text{productive of health}
\]

then:

\[
\text{beneficially} = \text{productively of health}
\]

and:

\[
\text{benefactor} = \text{producer of health}
\]

\subsection*{Teaching Point}

Definitions should remain stable across grammatical relatives.

\[
\boxed{
\text{A definition should survive its inflexions.}
}
\]

\section*{63. Definitions and Platonic Ideas}

Aristotle uses a rule against Platonists.

If a definition is meant to define a thing, it should also apply to the Idea if
one believes in Ideas.

\subsection*{Example}

If one defines man by adding:

\[
\text{mortal}
\]

the definition will not apply to the absolute Man, since the Platonic Idea is
not mortal.

\subsection*{Teaching Point}

A Platonist definition must fit the Form as well as the sensibles.

\[
\boxed{
\text{Definitions can be tested against an opponent's metaphysical commitments.}
}
\]

\section*{64. Ambiguous Terms Should Not Receive One Common Definition}

If a term is ambiguous, a single common definition will fail.

\subsection*{Example}

Life in plants and life in animals may not be the same in the relevant sense.

A definition that applies equally to both may fail to define either properly.

\[
\text{ambiguous term} \Rightarrow \text{multiple definitions}
\]

\subsection*{Teaching Point}

Synonymous terms can share one definition; ambiguous terms cannot.

\[
\boxed{
\text{One definition belongs to one meaning.}
}
\]

\section*{65. Handle Ambiguity Strategically}

Aristotle gives practical advice.

In questioning, treat ambiguous terms as if they were synonymous unless the
ambiguity has been secured. This may show that the answerer's definition fails
to apply to the full range.

In answering, distinguish the senses.

\subsection*{Whiteboard}

\[
\text{questioner: press the shared term}
\]

\[
\text{answerer: divide the senses}
\]

\subsection*{Teaching Point}

Dialectic requires both logical clarity and strategic awareness.

\[
\boxed{
\text{The questioner presses unity; the answerer distinguishes meanings.}
}
\]

\section*{66. Complex Terms: Test Part by Part}

If a complex term is defined, remove the definition of one element and see
whether the rest defines the remainder.

\subsection*{Example}

If finite straight line is defined, subtract the definition of finite line and
ask whether the rest defines straight.

\subsection*{Teaching Point}

A definition of a complex should distribute correctly across its elements.

\[
\boxed{
\text{Test a complex definition by subtracting its parts.}
}
\]

\section*{67. Do Not Merely Exchange Words}

A definition should not simply replace each term with another term.

It should explain by phrases, especially where the term is simple or obscure.

\[
\text{name swap} \neq \text{definition}
\]

\subsection*{Teaching Point}

A definition is not a synonym; it is an explanatory account.

\[
\boxed{
\text{Definition must unpack, not merely rename.}
}
\]

\section*{68. Equimembral Definitions Are Suspect}

If the definition has exactly the same number of elements as the thing defined,
it may simply be exchanging word for word rather than explaining.

\subsection*{Example}

\[
\text{white man} \rightarrow \text{pellucid mortal}
\]

This is less a definition than a re-labeling.

\subsection*{Teaching Point}

A definition should unpack, not merely rename.

\[
\boxed{
\text{Equimembral exchange is usually not explanation.}
}
\]

\section*{69. The Definition Must Describe Something Real if the Term Is Real}

If the term defined is real, but the definition describes something impossible,
the definition fails.

\subsection*{Example}

\[
\text{white} = \text{colour mingled with fire}
\]

If colour cannot be mingled with body in that way, the definition describes
something unreal, while white exists.

\subsection*{Teaching Point}

A definition of a real thing must not describe an impossibility.

\[
\boxed{
\text{The account must be as real as the thing accounted for.}
}
\]

\section*{70. Do Not Define a Thing Only in Its Good or Perfect Condition}

A definition should define the thing simply, not only its good or perfect
instance.

\subsection*{Examples}

A rhetorician should not be defined as:

\[
\text{one who can always see what will persuade and omit nothing}
\]

because this defines a good rhetorician, not a rhetorician simply.

A thief should not be defined only as one who successfully steals in secret if
the essence concerns the wish or choice to steal.

\subsection*{Teaching Point}

Do not define the thing by its perfected instance.

\[
\boxed{
\text{The good } S \neq \text{the definition of } S.
}
\]

\section*{71. Do Not Define What Is Desirable for Itself by What It Produces}

If something is desirable for its own sake, do not define it merely by what it
produces or preserves.

\subsection*{Bad Examples}

\[
\text{justice} = \text{what preserves the laws}
\]

\[
\text{wisdom} = \text{what produces happiness}
\]

These may name effects, but they do not state the intrinsic essence.

\subsection*{Teaching Point}

Intrinsic goods should be defined by what they are, not merely by what they do.

\[
\boxed{
\text{Do not reduce intrinsic good to instrumental effect.}
}
\]

\section*{72. Beware Definitions as ``A and B''}

Aristotle warns against defining something simply as:

\[
A \text{ and } B
\]

\subsection*{Example}

\[
\text{justice} = \text{temperance and courage}
\]

If two people each possess one of these, then together they might seem to have
justice, while neither individually is just.

This shows that the definition has confused whole and parts.

\subsection*{Teaching Point}

A thing is not usually identical with a mere list of components.

\[
\boxed{
\text{A whole is not merely its parts named side by side.}
}
\]

\section*{73. Beware Definitions as ``Product of A and B''}

If something is defined as the product of A and B, ask whether A and B can
naturally produce one thing.

Also ask whether the whole and parts belong in the same subject.

\[
\text{product of } A \text{ and } B
\Rightarrow
\text{explain how one thing results}
\]

\subsection*{Teaching Point}

A product-definition must explain the unity of the product.

\[
\boxed{
\text{Naming ingredients is not enough.}
}
\]

\section*{74. Whole and Parts Must Be Properly Related}

The whole should perish when the parts perish, but the parts need not perish
when the whole perishes.

If a proposed definition reverses this or fails to preserve whole-part
relations, it is suspect.

\subsection*{Teaching Point}

A compound definition must respect the logic of whole and part.

\[
\boxed{
\text{A whole depends on its parts, but the parts need not depend on the whole.}
}
\]

\section*{75. State the Manner of Composition}

If something is defined as a composition, it is not enough to name the elements.

One must state the manner of their composition.

\subsection*{Example}

The same elements may produce flesh or bone depending on how they are
compounded.

\[
\text{elements} + \text{mode of composition} = \text{compound}
\]

\subsection*{Teaching Point}

A compound is not merely its materials, but its materials arranged in a certain
way.

\[
\boxed{
\text{Composition requires order, not merely ingredients.}
}
\]

\section*{76. Attack a Part if the Whole Is Hard to Attack}

If one cannot attack the whole definition, attack a part of it.

If a part is destroyed, the definition is destroyed.

\[
\text{destroy part} \Rightarrow \text{destroy whole definition}
\]

\subsection*{Teaching Point}

A definition is only as strong as its weakest element.

\[
\boxed{
\text{When the whole is too hard, attack a part.}
}
\]

\section*{77. Clarify Obscure Definitions Before Attacking}

If a definition is obscure, Aristotle recommends reshaping or clarifying it in
order to find a handle for attack.

The answerer must either accept the clarified version or explain the intended
sense more clearly.

\subsection*{Teaching Point}

Clarification is itself a dialectical weapon.

\[
\boxed{
\text{Make the definition clear enough to be tested.}
}
\]

\section*{78. Propose a Better Definition}

Aristotle ends with a practical rule: propose a better definition.

If the new definition is clearer and more indicative of the object, the old one
is demolished, since there cannot be more than one definition of the same thing.

\[
\text{better definition} \Rightarrow \text{old definition defeated}
\]

\subsection*{Teaching Point}

The best attack on a bad definition is often a better definition.

\[
\boxed{
\text{Replace the bad account with a better account of the essence.}
}
\]

\section*{79. The Architecture of Book VI}

Book VI moves through the following arc:

\[
\text{what definition must do}
\rightarrow
\text{five tests of definition}
\rightarrow
\text{obscurity}
\rightarrow
\text{redundancy}
\]

\[
\rightarrow
\text{prior and more intelligible terms}
\rightarrow
\text{genus and differentiae}
\rightarrow
\text{relative definitions}
\rightarrow
\text{differentiae and essence}
\]

\[
\rightarrow
\text{cause/effect and time}
\rightarrow
\text{opposites and privations}
\rightarrow
\text{ambiguity}
\rightarrow
\text{complex terms}
\]

\[
\rightarrow
\text{whole and parts}
\rightarrow
\text{compound definitions}
\rightarrow
\text{better replacement definitions}
\]

\subsection*{Teaching Point}

Book VI teaches how to test whether a proposed definition truly states the
essence.

\[
\boxed{
\text{Definition is the dialectical test of essence.}
}
\]

\section*{80. Summary of Main Ideas}

\begin{enumerate}[itemsep=0.5em]
    \item Book VI gives commonplace rules for testing definitions.

    \item A definition is convertible with the subject and states the essence.

    \item The five tests ask whether the expression is true of the subject,
    places the object in the proper genus, is peculiar, states the essence, and
    is correctly formed.

    \item Some definition problems are handled by the rules for accident, genus,
    and property.

    \item Incorrect definitions are often obscure or redundant.

    \item Ambiguous, metaphorical, strange, or unfamiliar language makes a
    definition obscure.

    \item A good definition should help reveal its contrary.

    \item A definition should identify what it defines without requiring an
    external label.

    \item Redundant definitions include universal attributes, removable
    additions, repeated meanings, or particulars already included under a
    universal.

    \item A definition must use terms that are prior and more intelligible.

    \item Scientific definition should use what is more intelligible absolutely,
    even if teaching often begins with what is more intelligible to us.

    \item Genus and differentiae are prior to species.

    \item A definition should not explain the prior through the posterior.

    \item A definition should not define the definite through the indefinite.

    \item Opposites should not usually be defined through opposites, except in
    essentially relative cases.

    \item A definition must not contain the term defined, even indirectly.

    \item Co-ordinate members of a division should not be defined through one
    another.

    \item The superior should not be defined through the subordinate.

    \item A definition should place the object in its proper and nearest genus.

    \item Relative terms must be defined in relation to their proper correlates
    and ends.

    \item Differentiae must belong to the genus, have co-ordinate alternatives,
    and form species when added to the genus.

    \item A genus should not be divided merely by negation unless the object is
    privative.

    \item Species, genera, and differentiae must not be confused.

    \item Differentia signifies a quality of the genus and belongs essentially.

    \item Locality, affection, accident, and destructive states are usually not
    differentiae.

    \item Things should be defined by natural function, primary relation, and the
    proper subject of their states or affections.

    \item Definitions must not confuse cause and effect.

    \item Definition and defined term should match in temporal force and degree.

    \item Definitions must not define one essence by disconnected alternatives.

    \item A definition should remain coherent when its terms are unpacked.

    \item Privative terms must say both what is lacking and where it is naturally
    due.

    \item Ambiguous terms should not receive one common definition.

    \item Complex definitions should be tested part by part.

    \item A definition is not a mere synonym or word exchange.

    \item A definition of a real thing must not describe an impossibility.

    \item A thing should not be defined only in its good or perfect condition.

    \item Intrinsic goods should not be defined merely by what they produce.

    \item Compound definitions must explain the unity and manner of composition.

    \item If the whole definition is hard to attack, attack a part.

    \item A better definition can demolish a worse definition.
\end{enumerate}

\section*{Final Framing}

\[
\boxed{
\textit{Topics}, \text{ Book VI, teaches how to test whether a proposed definition truly states the essence.}
}
\]

Book VI teaches that a definition must be clear, non-redundant, non-circular,
prior, intelligible, properly placed in genus, properly divided by
differentiae, peculiar to the object, and expressive of essence.

A definition is not a metaphor, not a list, not a synonym, not a cause, not an
effect, and not a description of the best case. It is an account of what the
thing is.

\[
\boxed{
\text{A definition is the disciplined statement of essence.}
}
\]

\end{document}